\documentclass[11pt]{article}
\newcommand{\ReportShortTitle}{Initial ODMA decoding}
\newcommand{\ReportUpdated}{5 September 2026}
\usepackage[T1]{fontenc}
\usepackage[utf8]{inputenc}
\usepackage{lmodern}
\usepackage{microtype}
\usepackage[a4paper,margin=24mm,headheight=14pt]{geometry}
\usepackage{amsmath,amssymb,mathtools}
\usepackage{booktabs,tabularx,array}
\usepackage{enumitem}
\usepackage{xcolor}
\usepackage{fancyhdr}
\usepackage{titlesec}
\usepackage[hidelinks]{hyperref}

\definecolor{ink}{HTML}{18212B}
\definecolor{accent}{HTML}{135F73}
\definecolor{muted}{HTML}{5E6A73}
\definecolor{pale}{HTML}{EEF5F7}
\color{ink}
\hypersetup{colorlinks=true,linkcolor=accent,urlcolor=accent}
\setlength{\parindent}{0pt}
\setlength{\parskip}{0.55em}
\setlist{nosep,leftmargin=1.55em}
\setcounter{secnumdepth}{2}
\titleformat{\section}{\Large\bfseries\color{accent}}{\thesection}{0.6em}{}
\titleformat{\subsection}{\large\bfseries\color{ink}}{\thesubsection}{0.55em}{}
\pagestyle{fancy}
\fancyhf{}
\fancyhead[L]{\small ODMA--URA research record}
\fancyhead[R]{\small\color{muted}\ReportShortTitle}
\fancyfoot[C]{\thepage}
\renewcommand{\headrulewidth}{0.3pt}
\newcolumntype{Y}{>{\raggedright\arraybackslash}X}
\newcommand{\R}{\mathbb{R}}
\newcommand{\C}{\mathbb{C}}
\newcommand{\Z}{\mathbb{Z}}
\newcommand{\F}{\mathbb{F}}
\newcommand{\one}{\mathbf{1}}
\newcommand{\given}{\,|\,}
\newcommand{\CN}{\mathcal{CN}}
\DeclareMathOperator{\diag}{diag}
\DeclareMathOperator{\Cov}{Cov}
\DeclareMathOperator{\Var}{Var}
\DeclareMathOperator{\supp}{supp}
\DeclareMathOperator*{\argminop}{arg\,min}
\DeclareMathOperator*{\argmaxop}{arg\,max}
\newcommand{\evidence}[1]{\textcolor{accent}{\textbf{Evidence status:}} #1}
\providecommand{\ReportUpdated}{30 August 2026}
\newcommand{\ProjectReport}[4]{%
  \pagenumbering{roman}
  \begin{titlepage}
    \vspace*{18mm}
    {\large\color{accent}\textbf{ODMA--URA research record}}\\[12mm]
    {\huge\bfseries #1}\\[5mm]
    {\Large\color{muted}#2}\\[18mm]
    \begin{tabularx}{\textwidth}{@{}p{0.16\textwidth}Y@{}}
      Stage & #3\\
      Updated & \ReportUpdated\\
      Status & #4
    \end{tabularx}
    \vfill
    \textcolor{muted}{This report is a chronological research record. Measured results, interpretations, and open
    hypotheses are labelled separately. Detailed numerical ledgers remain in \texttt{results/}.}
  \end{titlepage}
  \pagenumbering{arabic}
}


\begin{document}
\ProjectReport{Initial ODMA-aware decoding}{The structured signal model, BP/EP derivation, and V1--V4 development}{1 of 5}
{Historical stage reconstructed from the original derivations and retained implementation}

\section{Position in the research story}

The project began with a receiver-side question: if every message waveform has explicit on--off division multiple
access (ODMA) structure, can a decoder exploit the resulting sparse factor graph more effectively than a generic
global sparse-recovery method? This report retains the algebra used to answer that question. It is intentionally more
detailed than a present-day summary because these equations explain how the early algorithms were constructed and why
their failure modes led to the later codebook-design programme.

The central receiver was a two-stage hybrid:
\begin{enumerate}
 \item a linear-Gaussian resource stage, which compresses interfering contributions to means and covariances;
 \item a discrete block stage, which enumerates valid local count vectors and uses the exact local codebook.
\end{enumerate}
The stages exchange extrinsic Gaussian information. Successive V1--V4 variants changed the observation model,
covariance approximation, schedule, and fading treatment; they did not change the underlying unsourced target.

\evidence{Historical reconstruction. The factorisation and current Graph-BP implementation are directly verifiable;
some superseded V1--V4 variants survive only as research records and should not be treated as current baselines.}

\section{ODMA--URA signal model}

\subsection{Patterns, local codebooks, and global messages}

Let $n$ be the number of channel uses and let $Q$ be the number of ODMA placement patterns. Pattern $q$ occupies
$d_q$ resources $S_q\subseteq\{1,\ldots,n\}$ and is represented by
\[
 P_q\in\{0,1\}^{n\times d_q},\qquad P_q^HP_q=I_{d_q}.
\]
If $x_q\in\C^{d_q}$ is a local signal, $P_qx_q$ embeds it into the full resource grid. Pattern $q$ owns $V_q$ local
codewords, stored as rows of
\[
 C_q\in\C^{V_q\times d_q},\qquad c_{q,v}=C_q[v,:]^T.
\]
The global message $(q,v)$ transmits
\[
 \phi_{q,v}=P_qc_{q,v}\in\C^n.
\]
When every pair is legal, the total message alphabet is $M=\sum_qV_q$ and the explicit global codebook is
\[
 \Phi=[P_1C_1^T\mid P_2C_2^T\mid\cdots\mid P_QC_Q^T]\in\C^{n\times M}.
\]
Column normalisation $\|\phi_m\|_2^2=1$ fixes the transmitted energy of every message.

The $B$ payload bits select one of $M=2^B$ columns. They are labels, not binary channel symbols: the entries of
$\phi_m$ may be real or complex amplitudes, so each physical channel use may carry magnitude and phase.

\subsection{Unsourced count target}

For $K_a$ active devices, define
\[
 a_q\in\Z_+^{V_q},\qquad a_{q,v}=\text{number of devices that chose }(q,v).
\]
The local superposition in block $q$ and the global resource signal are
\[
 x_q=C_q^Ta_q\in\C^{d_q},\qquad
 s=\sum_{q=1}^QP_qx_q=\Phi a,qquad \one^Ta=K_a.
\]
The decoder seeks the multiset $a$, not device identities. If two devices send the same message then $a_m=2$; this is
a valid message collision, not two separately identifiable users.

\subsection{Observation models used during V1--V4}

The original sequence can be written as increasingly difficult versions of the same resource signal.

\paragraph{V1: scalar, known unit channel.}
\[
 y=s+z=\sum_qP_qC_q^Ta_q+z,
 \qquad z\sim\CN(0,\sigma^2I_n).
\]

\paragraph{V2: multiple antennas with a known common signature.}
For $M_{\rm ant}$ receive antennas and known $h\in\C^{M_{\rm ant}}$,
\[
 Y=sh^T+Z,
 \qquad Z_{i,:}\sim\CN(0,\sigma^2I).
\]
Let $\gamma=\|h\|_2^2$. Matched combining gives the sufficient scalar observation
\[
 \bar y=\frac{Yh^*}{\gamma}=s+\bar z,
 \qquad \bar z\sim\CN(0,\sigma^2\gamma^{-1}I_n).
\]
In the implementation $h=\one$, so $\gamma=M_{\rm ant}$. The orthogonal antenna projector
\[
 Q_\perp=I-\frac{hh^H}{\gamma}
\]
contains no signal under this model and can estimate noise because $Q_\perp Y_i=Q_\perp Z_i$.

\paragraph{V3: message-dependent unknown fading.}
With channel vector $h_u$ for device $u$ and selected message $m_u$,
\[
 Y=\sum_{u=1}^{K_a}\phi_{m_u}h_u^T+Z.
\]
Now activity and channel are bilinearly coupled. A count $a_m$ alone is no longer sufficient unless the fading
coefficients are known, integrated out, or replaced by an energy/support statistic.

\paragraph{V4: exact small-state reference.}
At tiny dimensions, the intended posterior over discrete counts and any retained channel state can be enumerated or
integrated exactly. This is a correctness reference, not a scalable receiver.

\section{Factor graph in block variables}

The early decoder used one variable node per ODMA block signal $x_q$ and one factor per resource. If coordinate $j$ of
block $q$ is placed at resource $i$, then $P_q[i,j]=1$. Define the neighbours
\[
 \mathcal N(i)=\{(q,j):P_q[i,j]=1\},\qquad
 \mathcal N(q)=S_q.
\]
Under V2, resource $i$ obeys
\[
 Y_i=h\sum_{(q,j)\in\mathcal N(i)}x_{q,j}+Z_i.
\]
The exact posterior factors into resource likelihoods and discrete block priors, but shared resources create loops.
The algorithm therefore approximates every incoming block-coordinate belief by a Gaussian.

\subsection{Stage 1: linear-Gaussian resource update}

Suppose edge $(q,j)\to i$ carries
\[
 x_{q,j}\sim\CN(m_{qj\to i},v_{qj\to i}).
\]
Excluding the candidate edge, the other contributions are approximated by
\[
 \mu_{i\setminus qj}=\sum_{(b,k)\in\mathcal N(i)\setminus(q,j)}m_{bk\to i},
 \qquad
 \nu_{i\setminus qj}=\sum_{(b,k)\in\mathcal N(i)\setminus(q,j)}v_{bk\to i}.
\]
After matched combining, the extrinsic scalar observation is
\[
 r_{i\to qj}=\bar y_i-\mu_{i\setminus qj}=x_{q,j}+\epsilon_{i\to qj},
 \qquad
 \Var(\epsilon_{i\to qj})=\nu_{i\setminus qj}+\sigma^2/\gamma.
\]
For vector block updates, stacking the relevant coordinates gives a Gaussian cavity
\[
 q^{\setminus q}(x_q)=\CN(x_q;r_q^{\rm cav},V_q^{\rm cav}).
\]

\subsection{Stage 2: exact discrete block update}

The structural constraint is
\[
 x_q=C_q^Ta_q,qquad a_q\in\mathcal A_q\subset\Z_+^{V_q}.
\]
For a Poisson approximation with rate $\lambda$ per message,
\[
 p(a_q)=\prod_{v=1}^{V_q}e^{-\lambda}\frac{\lambda^{a_{q,v}}}{a_{q,v}!}.
\]
Given the Gaussian cavity, the posterior weight of state $a\in\mathcal A_q$ is
\[
 \omega_q(a)\propto p(a)
 \exp\!\left[-(C_q^Ta-r_q^{\rm cav})^H(V_q^{\rm cav})^{-1}
 (C_q^Ta-r_q^{\rm cav})\right].
\]
Normalising these weights yields the exact moments inside the enumerated local state set:
\begin{align*}
 \widehat a_q&=\sum_{a\in\mathcal A_q}\omega_q(a)a,\\
 \widehat x_q&=\sum_{a\in\mathcal A_q}\omega_q(a)C_q^Ta,\\
 \widehat\Sigma_q&=\sum_{a\in\mathcal A_q}\omega_q(a)
 (C_q^Ta-\widehat x_q)(C_q^Ta-\widehat x_q)^H.
\end{align*}
The block MAP state is $a_q^{\rm MAP}=\argmaxop_{a\in\mathcal A_q}\omega_q(a)$.

The state count explains the cost. If at most $k_{\max}$ local messages are nonzero and each count is at most
$c_{\max}$, then
\[
 |\mathcal A_q|=\sum_{k=0}^{k_{\max}}\binom{V_q}{k}c_{\max}^{k}.
\]
Truncation makes the decoder executable, but it also changes the inference domain.

\subsection{Extrinsic EP site update}

Write a Gaussian in natural parameters $(\Lambda,\eta)=(\Sigma^{-1},\Sigma^{-1}\mu)$. EP divides the moment-matched
posterior by the cavity:
\[
 \Lambda_q^{\rm site}=\widehat\Sigma_q^{-1}-(V_q^{\rm cav})^{-1},
 \qquad
 \eta_q^{\rm site}=\widehat\Sigma_q^{-1}\widehat x_q-(V_q^{\rm cav})^{-1}r_q^{\rm cav}.
\]
These site parameters are damped and returned to the linear system. The current implementation forms the cavity by
omitting the old block site from the global precision system rather than subtracting nearly equal posterior and site
parameters after solving.

\subsection{EM updates used by the approximation}

If $M=\sum_qV_q$, the Poisson-rate update is
\[
 \lambda^{+}=\frac{1}{M}\sum_q\one^T\widehat a_q.
\]
Let $\widehat y=\sum_qP_q\widehat x_q$. A moment-corrected noise update has the form
\[
 (\sigma^2)^{+}=\frac{\|Y-\widehat yh^T\|_F^2+gamma\sum_q\operatorname{tr}(\widehat\Sigma_q)}
 {nM_{\rm ant}}.
\]
These are empirical-Bayes updates inside an approximate posterior, not ground-truth parameter estimates.

\section{Implementation sequence and what each change meant}

\begin{tabularx}{\textwidth}{@{}lYY@{}}
\toprule
Variant & Structural change & Intended diagnostic \\
\midrule
V1 & Scalar, no-fading resource messages & Establish the two-stage algebra in the simplest channel \\
V2 & Common known multi-antenna signature & Test antenna combining without bilinear channel uncertainty \\
V2b & Block/full covariance rather than diagonal moments & Measure the value of within-block correlations \\
V2c & Onsager-like self-response correction & Reduce feedback through loopy shared resources \\
V2d & Serial/turbo block schedule & Use freshly updated cavities rather than stale parallel messages \\
V2e & Global EP with cavity-by-omission & Avoid unstable subtraction when forming cavities \\
V2f & Fractional EP site power & Limit information injection when sites become overconfident \\
V3 & Explicit latent fading variables & Represent the physical bilinear channel rather than pre-equalise it \\
V3a & Energy/support scoring inspired by noncoherent aggregation & Retain multi-antenna support evidence at lower cost \\
V4 & Exact small-state likelihood/reference & Separate approximation failure from model/code errors \\
\bottomrule
\end{tabularx}

Fractional EP replaces a full site update by a powered one. In natural coordinates a simple damped form is
\[
 \Lambda_q^+ = (1-\alpha)\Lambda_q^{\rm old}+\alpha\beta\Lambda_q^{\rm site},
 \qquad
 \eta_q^+ = (1-\alpha)\eta_q^{\rm old}+\alpha\beta\eta_q^{\rm site},
 \qquad 0<\beta\le1.
\]
It can stabilise an iteration by shrinking the update, but does not correct an inaccurate local likelihood or remove
short loops. Similarly, a serial schedule changes dynamics, not the fixed-point objective.

\section{Why the graph decoder became unstable}

The observed high-load failure has a concrete EP mechanism. When the discrete posterior places nearly all mass on one
state, $\widehat\Sigma_q$ becomes very small. Then $\widehat\Sigma_q^{-1}$ and hence the site precision become large.
That strong precision tightens the cavities of every overlapping block. With many short loops, the next blocks become
overconfident in turn, and the feedback gain can exceed one. Damping and fractional powers slow this amplification;
they do not guarantee that the Gaussian projection is accurate.

Four additional limitations remained:
\begin{enumerate}
 \item the independent Poisson prior does not impose the exact multinomial constraint $\one^Ta=K_a$;
 \item local state caps may exclude the true high-occupancy block state;
 \item ODMA reuses a small number of support subspaces, creating correlated columns and local ambiguities;
 \item under unknown fading, activity and channel scale/phase are not separately identified without extra structure.
\end{enumerate}

\section{Baseline comparison and historical conclusion}

The important baselines were global OMP/NNOMP, linear MMSE, SIC, AMP-like iterations, local BlockMAP, and exact tiny
references. Their assumptions differ. In particular, NNOMP was usually supplied the true $K_a$ and therefore received
oracle knowledge of how much total mass to select. Exact BlockMAP ignored or only approximated cross-block
interference. V4 was too small to be a practical competitor.

Qualitatively, the structured graph decoder could work in easy regimes but did not show a stable advantage as load,
overlap, or fading uncertainty increased. Global NNOMP was simpler and usually stronger in the comparisons that drove
the project. This did not establish that ODMA was information-theoretically inferior. It established that adding more
local message-passing machinery was not isolating the dominant loss.

The next question was therefore causal: once the true active columns are supplied, are ODMA columns intrinsically hard
to fit, or is the main loss the search for those columns? Report 2 places every construction in one global count model,
catalogues the decoder objectives tried, and answers that question with an oracle-support control.

\end{document}
